


\section{Equilibrio estático}

\subsection{Primera condición de equilibrio}
La primera condición de equilibrio establece que un cuerpo se encuentra en equilibrio traslacional si la suma vectorial de todas las fuerzas que actúan sobre él es igual a cero.
\begin{equation}
    \sum \vec{F} = 0 \quad \Rightarrow \quad \sum F_x = 0, \quad \sum F_y = 0, \quad \sum F_z = 0
\end{equation}

\subsection{Segunda condición de equilibrio}
La segunda condición de equilibrio se refiere al equilibrio rotacional. Establece que la suma de todos los momentos de torsión (torcas) que actúan sobre un cuerpo debe ser cero.
\begin{equation}
    \sum \vec{\tau} = 0
\end{equation}

\section{Momentos de torsión (torca)}

El momento de torsión o torca es una medida de la tendencia de una fuerza a hacer rotar un objeto alrededor de un eje. Se define como:
\begin{equation}
    \vec{\tau} = \vec{r} \times \vec{F} \quad \Rightarrow \quad \tau = r \cdot F \cdot \sin(\theta)
\end{equation}
donde:
\begin{itemize}
    \item $ \tau $ es la torca,
    \item $ r $ es la distancia perpendicular desde el eje de rotación hasta la línea de acción de la fuerza,
    \item $ F $ es la magnitud de la fuerza,
    \item $ \theta $ es el ángulo entre el vector de posición y la fuerza.
\end{itemize}

\subsection{Signo de la torca}
\begin{itemize}
    \item \textbf{Torca positiva:} tiende a girar el cuerpo en sentido antihorario.
    \item \textbf{Torca negativa:} tiende a girar el cuerpo en sentido horario.
\end{itemize}

\section{Centroides}

El centroide es el punto donde se considera que se concentra toda el área o el volumen de un cuerpo. En problemas de estática, el centroide es importante para determinar el centro de gravedad.

\subsection{Coordenadas del centroide para figuras planas}
\begin{align}
    \bar{x} &= \frac{\int x \, dA}{\int dA} = \frac{1}{A} \int x \, dA \\
    \bar{y} &= \frac{\int y \, dA}{\int dA} = \frac{1}{A} \int y \, dA
\end{align}

\subsection{Centroide de figuras compuestas}
Para figuras compuestas, el centroide se calcula como el promedio ponderado de los centroides de cada parte:
\begin{align}
    \bar{x} &= \frac{\sum (A_i \cdot \bar{x}_i)}{\sum A_i} \\
    \bar{y} &= \frac{\sum (A_i \cdot \bar{y}_i)}{\sum A_i}
\end{align}

\begin{table}[h!]
    \centering
    \begin{tabular}{|c|c|c|}
        \hline
        \textbf{Figura} & \textbf{Área} & \textbf{Centroide} \\
        \hline
        Rectángulo & $A = b \cdot h$ & $\bar{x} = b/2, \bar{y} = h/2$ \\
        Triángulo & $A = \frac{1}{2} b \cdot h$ & $\bar{x} = b/3, \bar{y} = h/3$ (desde el vértice) \\
        Círculo & $A = \pi r^2$ & $\bar{x} = 0, \bar{y} = 0$ (centro) \\
        \hline
    \end{tabular}
    \caption{Centroides de figuras comunes}
\end{table}

\section{Ejercicios propuestos}

\begin{enumerate}
    \item Una viga de 6 metros está sometida a fuerzas en sus extremos. Calcula la torca neta y determina si está en equilibrio rotacional.
    \item Encuentra el centroide de una figura compuesta formada por un rectángulo de $4 \times 6$ cm y un triángulo encima.
\end{enumerate}